\section{The original signed return from the two universal Hankel ratios}

Retain CTR21--54's complete original quartet, $k=4l+1$,
$l\geq1$, $m\geq1$, $e=1+k(m-1)$, and
$q=e(k+1)^2$. In particular the two high relation ranks
are $q$ and $q+1$. The original real source density is
$\sigma(y)=|\Gamma(1/4+iy/2)|^2/(2\pi)$ with full mass
$c_\sigma=\sqrt{2\pi}$. Use the proved finite threshold
$k\geq2048\sqrt{d+g}$, where $d=\delta^2$ and $g=\gamma^2$
have their original domains, and set $\epsilon=2^{-10}$.
The following constants are exactly the ones in CTR23,
CTR38 and CTR45:
\[
\begin{gathered}
 a_k=\frac{k(k+2)(d-g)}{3q},\quad
 \tau_\epsilon=\frac{k^2(d+g)}{q^2\epsilon^2},\quad
 \delta_\chi=\frac{q\tau_\epsilon^2}{2(1-\tau_\epsilon)},\\
 E=E_{\chi,l}=l\{|a_k|(\epsilon^{-2}-64^{-2})+2\delta_\chi\},\quad
 \Delta_f=\frac{l(16l^2-12l-1)}{12q^2\epsilon^2},\\
 \kappa_{I,r}=\frac{2\epsilon r^2}{\sqrt{\cos1}}
     \left(\frac43\,2^{-13}e^{2\pi}\right)^q,\quad
 \kappa_F=\frac{512q}{59\sqrt{\cos1}}e^{-19q},\quad
 L_r=r\log(1+\kappa_{I,r}+\kappa_F).
\end{gathered}\tag{HSR1}
\]
Here $a_k$ is CTR's signed quadratic-density coefficient;
the separate arithmetic deficit sums below have superscripts
$+$ and $-$ and are not this coefficient.

For $r=q,q+1$, put
\[
 (\mathsf M_{a,r})_{ij}=\int_{\mathbb R}y^{2a+i+j}\sigma(y)\,dy,
 \qquad
 h_r=\log\frac{\det\mathsf M_{q+l,r}}{\det\mathsf M_{q,r}},
 \quad 0\leq i,j<r.
\]
These are exactly CTR34's original full-source forms. Let
\[
\begin{gathered}
 C_{\rm rel}=\log D_{l,q-1}+\log D_{l,q}-\log\mathfrak r_\chi
                         -\log D_{l,2q-1}-\log D_{l,2q},\\
 D_{l,N}=\prod_{a=0}^{N}\prod_{j=0}^{2l-1}(a+j+\tfrac12),\qquad
 \mathfrak r_\chi=\frac{\int|f(c+iy)\chi(c+iy)|^2\sigma(y)\,dy}
                         {\int|\chi(c+iy)|^2\sigma(y)\,dy}>0,\\
 H_k^*=-C_{\rm rel}-h_q-h_{q+1},\qquad
 B_k^-=2E+(2q+1)\Delta_f+2(L_q+L_{q+1}),\quad
 B_k^+=2E+2(L_q+L_{q+1}).
\end{gathered}\tag{HSR2}
\]
The original $f$, full polynomial $\chi$, variable $c+iy$ and
all masses are retained in the low-endpoint ratio. JSR24b
already provides its finite coefficient-Gram expression.

The complete signed Gamma return satisfies
\[
 \boxed{\quad H_k^*-B_k^-\leq
 \mathcal H_k:=\mathcal R_k^0-\mathcal R_k^\sigma
             =-\mathfrak T_k
 \leq H_k^*+B_k^+.\quad}\tag{HSR3}
\]
To prove this, retain CTR47's full error
$\varepsilon_r=\log\mathfrak R_{q+r-1}-h_r$.
Its actual four-tail decomposition is
\[
 \varepsilon_r=d_\chi+d_f
   +t_{H_{\chi,f}}-t_{H_{\chi,0}}
   -t_{H_{\mathrm{pow},l}}+t_{H_{\mathrm{pow},0}},
\]
where the bulk proof gives $|d_\chi|\leq E$ and
$0\leq d_f\leq r\Delta_f$, and each of the four
nonnegative tail logarithms lies in $[0,L_r]$.
Consequently
$-E-2L_r\leq\varepsilon_r\leq E+r\Delta_f+2L_r$.
Adding at the two distinct ranks $q,q+1$ gives
$-B_k^+\leq\varepsilon_q+\varepsilon_{q+1}\leq B_k^-$.
The exact GSR/BRD/CTR identity is
$\mathfrak T_k=C_{\rm rel}+h_q+h_{q+1}
                 +\varepsilon_q+\varepsilon_{q+1}$.
Negating that entire equality proves HSR3. In particular
the one-sided multiplier error enters only $B_k^-$.
The exact finite interval width is
\[
 B_k^-+B_k^+=4E+(2q+1)\Delta_f+4(L_q+L_{q+1})=o(q).
 \tag{HSR4}
\]
Indeed CTR31--33 give $E=O(k/e)$, and
$(2q+1)\Delta_f\leq k^3/(4q\epsilon^2)=O(k/e)$ using
$2q+1\leq4q$ and $\sum_{j<l}(2j+1/2)^2\leq k^3/16$.
Each $L_r$ has the explicit exponentially decreasing bound
in HSR1. Since $q=e(k+1)^2$, these estimates give $O(k)=o(q)$
for $m=1$, and $O(1)=o(q)$ for each fixed $m\geq2$.
This error estimate leaves the value of $H_k^*$ intact.

Keep the original arithmetic deficit sums and their accepted
degree-specific outer bounds
\[
 a_k^-=\Xi_{k,q-1}+\Xi_{k,q}\leq D_k^-,\quad
 a_k^+=\Xi_{k,2q-1}+\Xi_{k,2q}\leq D_k^+,\quad
 -a_k^-\leq\delta_k^\sigma\leq a_k^+,
 \tag{HSR5}
\]
where $D_k^\pm$ are precisely the finite quantities in
HMR44/JSR19--21, retaining the original $q$. Addition to
HSR3 proves the actual receiving intervals
\[
\begin{aligned}
 H_k^*-B_k^--a_k^-&\leq\Delta_k^\Gamma
                     \leq H_k^*+B_k^++a_k^+,\\
 \mathcal B_k^0+H_k^*-B_k^--a_k^-&\leq\mathcal B_k^{\rm ar}
                     \leq\mathcal B_k^0+H_k^*+B_k^++a_k^+.
\end{aligned}\tag{HSR6}
\]
The identities used are the unchanged
$\Delta_k^\Gamma=\mathcal H_k+\delta_k^\sigma$ and
$\mathcal B_k^{\rm ar}=\mathcal B_k^0+\Delta_k^\Gamma$.
Replacing $a_k^\pm$ by $D_k^\pm$ gives explicit outer
intervals. For example the complete original HC residual obeys
\[
\begin{gathered}
 \max\{0,\mathcal B_k^0+H_k^*-B_k^--D_k^-
                              -4q\log(D_hk)\}
 \leq\mathcal R_k\\
 \leq\mathcal B_k^0+H_k^*+B_k^++D_k^+
                              -4q\log(D_hk).
\end{gathered}\tag{HSR7}
\]
This subtracts the same original threshold and uses its
proved residual nonnegativity. The narrower version using
$a_k^\pm$ follows identically. The original TWA coarse bounds
can also be retained by replacing $a_k^-$ with
$\min\{a_k^-,qE_k^+\}$ and $a_k^+$ with
$\min\{a_k^+,qE_k^-\}$; each minimum is a bound on the
same signed arithmetic correction, not a new observation.

For fixed $m\geq2$, the width of the explicit $D_k^\pm$
interval in HSR6 divided by $q\log k$ tends to
$192m+735$. Its width is exactly
$B_k^-+B_k^++D_k^-+D_k^+$; HSR4 divided by $q\log k$
tends to zero, and the two original ODD limits add to the
stated constant. The version with actual $a_k^\pm$ has
the corresponding upper limit at most $192m+735$.
For $m=1$, the retained arithmetic rate is the existing
$O_h((\log k/k)^{5/7})$ on the $kq$ scale.

\subsection{An exact parity factorization at ranks $q$ and $q+1$}

Let $\nu_j=\mu_{2j}/c_\sigma$ denote the positive rational
even moments from CTR50--53, and define
\[
 (\mathsf V_{a,n})_{ij}=\nu_{a+i+j},\quad 0\leq i,j<n,
 \qquad v_{a,n}=\det\mathsf V_{a,n},\quad v_{a,0}=1.
\]
If $n_+=\lceil r/2\rceil$ and $n_-=\lfloor r/2\rfloor$,
the permutation listing the even coefficient indices before
the odd ones gives the exact congruence
\[
 \Pi_r^*\mathsf M_{a,r}\Pi_r
 =c_\sigma\begin{pmatrix}
      \mathsf V_{a,n_+}&0\\0&\mathsf V_{a+1,n_-}
                    \end{pmatrix},\quad
 \det\mathsf M_{a,r}=c_\sigma^r v_{a,n_+}v_{a+1,n_-}.
 \tag{HSR8}
\]
The cross block vanishes because its exponent is odd and the
original $\sigma$ is even. The two diagonal entries follow
by putting the original indices equal to $2i,2j$, or
$2i+1,2j+1$. This proves the congruence without assigning
it to an independent new measure. Each block is positive
definite: its quadratic form integrates respectively
$y^{2a}|P(y^2)|^2$ or $y^{2a+2}|P(y^2)|^2$ against
$\sigma/c_\sigma$, and a nonzero polynomial cannot vanish
on a real interval. Thus all displayed determinants are
strictly positive rational numbers. The full mass factor
has been displayed before any cancellation.

The original $q=e(4l+2)^2$ is even; put $n=q/2$.
Using ranks $q=2n$ and $q+1=2n+1$ separately in HSR8 gives
\[
 \prod_{r=q}^{q+1}\frac{\det\mathsf M_{q+l,r}}
                             {\det\mathsf M_{q,r}}
 =\frac{v_{q+l,n}\,v_{q+l,n+1}\,v_{q+l+1,n}^{,2}}
        {v_{q,n}\,v_{q,n+1}\,v_{q+1,n}^{,2}}.
 \tag{HSR9}
\]
Therefore $H_k^*$ in HSR2 is exactly minus $C_{\rm rel}$
minus the logarithm of the positive rational number in
HSR9. Both original relation ranks remain present through
the separate sizes $n$ and $n+1$. This factorization is a
finite evaluation dictionary, not an asymptotic assignment
to that rational determinant ratio.
